
This paper reports two findings. The first concerns models: in every model $\times$ dataset cell tested, intermediate layers of a language model scored its own hallucinations better than its output layer did, using only the uncertainty statistics of the token distributions the model already computes en route to its answer. The second concerns measurement: the documented failure baseline this study was designed to compare against did not survive a provenance audit. Six protocol differences --- sample selection, prompt template, label rule, signal definition, numeric precision, and evaluation split --- moved the same cell's final-layer entropy AUROC from 0.5186 to 0.5928 before any new claim had been tested. We treat the second finding as a result in its own right.

The underlying problem is familiar. Open-weight LLMs hallucinate, and the cheapest detection signals --- confidence statistics read from the output distribution --- are architecture-fragile. In an archived record from an earlier study in this research line (hereafter the \textit{motivating record}), final-layer mean token entropy reached AUROC 0.66 on LLaMA-3-8B-Instruct but approximately 0.52, with inverted score direction, on base LLaMA-2-7B: the same statistic, on the same datasets. We cite these numbers only as motivating context; the audit above shows that they are bound to that record's protocol, and they are never used as same-protocol baselines here. The fragility itself, however, is the point. A confidence score whose validity flips with the checkpoint cannot be deployed reliably, and the standard alternatives are expensive: multi-sample semantic consistency costs roughly ten forward passes per query \citep{Farquhar2024}, and supervised hidden-state probes require labeled data and per-model training \citep{Azaria2023}.

At the same time, the field has largely converged on a diagnosis without adopting the least expensive remedy. That intermediate layers carry stronger truthfulness signal than the final layer is close to consensus [Orgad et al., 2024; Wang et al., 2026]. Yet the methods that exploit this observation reintroduce additional machinery: trained probes over hidden states [Azaria \& Mitchell, 2023; Suresh et al., 2025], geometric layer-selection criteria feeding supervised classifiers \citep{Wang2026}, or repeated sampling \citep{Zhang2024SLED}.

This leaves a specific gap. The raw logit lens --- projecting each layer's hidden state through the model's own unembedding --- yields a full next-token distribution at every depth, and the uncertainty statistics of those distributions have not previously been evaluated as hallucination-detection scores. Adjacent-layer KL divergence, a natural measure of how much the model revises its predicted distribution between layers, has been used only to steer decoding [Chuang et al., 2023; Wu et al., 2025], not scored for detection. The gap appears to have persisted partly through a community split --- interpretability work studies these statistics as computation signatures \citep{Ali2025}, while detection work defaults to probes and sampling --- and partly through an adjacent negative result: the nearest published measurement tracked final-prediction-token probability trajectories across layers and found certain and uncertain trajectories largely aligned \citep{Kim2025}, a result for a different statistic that could be over-read as covering this one.

The working premise is that the gap conceals signal because of where the standard signal is read. The final layer is not necessarily where the information disappears; it is where output calibration --- a tokenizer- and tuning-dependent presentation step --- reshapes it. Under this reading, when a model produces a correct answer, next-token candidate competition collapses at intermediate depth and remains collapsed; when it hallucinates, the decoded distribution stays wide and continues to shift until the output head consolidates it. The design principle that follows is to read before calibration: compute entropy, max-token probability, and adjacent-layer KL from every layer's logit-lens distribution --- obtained from one greedy generation plus one teacher-forced re-forward, with no sampling and no training --- and select the readout layer, signal, and score direction per model on a held-out selection split, treating calibration depth as a property of the checkpoint.

We test this premise at proof-of-concept scale under a staged verification design whose first tier asks only whether the signal exists. Across LLaMA-2-7B, Mistral-7B-v0.1, and LLaMA-3-8B-Instruct on TriviaQA and TruthfulQA --- 5,451 scored generations covering all six model $\times$ dataset cells --- at least one screened intermediate layer is class-separable in every cell, with selection-split corrected AUROC between 0.6092 and 0.7011 against a 0.55 existence gate, and the best screened intermediate layer beats the same sweep's final-layer entropy readout in every cell (point-estimate margins +0.0035 to +0.130). These are selection-split point estimates from a single seed, with no confidence intervals: the numbers that chose the (layer, signal, direction) tuples also grade them, and they must not be read as held-out performance. The locked test splits exist and were never read --- by any analysis, or by any decision made in writing this paper --- which preserves the deferred evaluation as an uncontaminated, pre-registered confirmation. Evaluating the frozen tuples there, along with cross-dataset transfer and KL--entropy fusion, is the explicitly unmeasured next tier of the staged design. Within this scope, the depth advantage is universal in direction and consistent with the calibration-suppression reading, though it does not establish that mechanism.

Four contributions follow. First, to our knowledge, the first AUROC evaluation of raw per-layer logit-lens uncertainty statistics as training-free, sampling-free hallucination scores: class-separability holds in all six cells at the existence tier. Second, a universal point-estimate depth advantage: the best screened intermediate layer --- concentrated at layers 28--31 --- exceeds the within-sweep final-layer baseline in every cell, extending the intermediate-over-final consensus to statistics that require no probe. Third, adjacent-layer KL as a detection signal in its own right: previously a decoding heuristic, it is the best signal in two of the six cells --- an even three-way split with entropy and max-probability --- and provides the top three intermediate scores in the binding LLaMA-2/TriviaQA cell. Fourth, a methodological contribution arising from a failure: protocol-internal validity anchoring. Having shown quantitatively that AUROC magnitudes for weak signals do not survive protocol changes --- direction transfers, magnitude does not --- we replace cross-run numeric anchors with identity-verified cache reuse, within-sweep baselines, and descriptive-only cross-run comparisons.
